% GENERATED FILE. Edit canonical lessons, then run npm run generate:books.
% canonical-source-hash: fnv1a64:032e70f02b3ba2d7
% canonical-lessons: HI-C543-vacan, HI-C543-qasam, HI-C543-razi, HI-C543-sapath, HI-C543-paband

\chapter{Keeping One's Word}
\label{ch:hi-keeping-one-s-word}
\hlchaptermodality{543}

\begin{chapteropening}
\textbf{By the end of this chapter:} \emph{I can talk about promises, oaths and being bound by them.}

Complete the last lesson of chapter 543: I can talk about promises, oaths and being bound by them.
\end{chapteropening}

\section[\texorpdfstring{\hi{वचन} (vacan)}{वचन (vacan)}]{\hi{वचन} (vacan) — a word given; a pledge}

\label{lesson:HI-C543-vacan}



\begin{quote}
Before the new one: say the Hindi for an offer, then the Hindi for a mediator.
\end{quote}

\subsection*{You'll want to know: \hi{वचन}}
\textbf{\hi{वचन}} — \emph{vacan} — \textquotedblleft{}a word given; a pledge\textquotedblright{}.

\begin{quote}
\hi{राजा} \hi{ने} \hi{रानी} \hi{को} \hi{एक} \hi{वचन} \hi{दिया।}
\end{quote}

\begin{quote}
\emph{Rājā ne rānī ko ek vacan diyā.}
\end{quote}

\begin{quote}
The king gave the queen his word.
\end{quote}

Masculine.

\begin{cousinweb}
\textbf{\hi{वचन}} (\emph{vacan}) is a tatsama, taken whole from Sanskrit \emph{vacana}, \textquotedblleft{}speaking, an utterance, a statement\textquotedblright{} (Monier-Williams); Hindi uses it for the word one gives.
\end{cousinweb}

\begin{grammarlens}[title={What you've built}]
One more word for saying what happens if.
\end{grammarlens}

\subsection*{Your turn}
\begin{itemize}
\raggedright
  \item \emph{Say it:} \emph{vacan}
  \item \emph{Say it:} \emph{vacan}, once more
  \item \emph{Say it:} \emph{Rājā ne rānī ko ek vacan diyā.} — The king gave the queen his word.
  \item \emph{Recall:} say the Hindi for an offer, then the Hindi for a mediator, then say \emph{vacan} again
\end{itemize}

\begin{tcolorbox}[breakable,colback=teal!4,colframe=teal!35!black,title={Before you move on}]
What does \hi{वचन} mean? (\textquotedblleft{}A word given; a pledge\textquotedblright{}.) Say it once more.
\end{tcolorbox}
\section[\texorpdfstring{\hi{क़सम} (qasam)}{क़सम (qasam)}]{\hi{क़सम} (qasam) — an oath; a vow}

\label{lesson:HI-C543-qasam}



\begin{quote}
Before the new one: say the Hindi for a mediator, then the Hindi for a word given.
\end{quote}

\subsection*{You'll want to know: \hi{क़सम}}
\textbf{\hi{क़सम}} — \emph{qasam} — \textquotedblleft{}an oath; a vow\textquotedblright{}.

\begin{quote}
\hi{मीरा} \hi{ने} \hi{क़सम} \hi{खाई} \hi{कि} \hi{वह} \hi{फिर} \hi{कभी} \hi{देर} \hi{से} \hi{नहीं} \hi{आएगी।}
\end{quote}

\begin{quote}
\emph{Mīrā ne qasam khāī ki vah phir kabhī der se nahīṃ āegī.}
\end{quote}

\begin{quote}
Mira swore that she would never be late again.
\end{quote}

Feminine.

\begin{cousinweb}
\textbf{\hi{क़सम}} (\emph{qasam}) is usually traced to Arabic \emph{qasam}, \textquotedblleft{}an oath\textquotedblright{}, through Persian.
\end{cousinweb}

\begin{grammarlens}[title={What you've built}]
One more word for saying what happens if.
\end{grammarlens}

\subsection*{Your turn}
\begin{itemize}
\raggedright
  \item \emph{Say it:} \emph{qasam}
  \item \emph{Say it:} \emph{qasam}, once more
  \item \emph{Say it:} \emph{Mīrā ne qasam khāī ki vah phir kabhī der se nahīṃ āegī.} — Mira swore that she would never be late again.
  \item \emph{Recall:} say the Hindi for a mediator, then the Hindi for a word given, then say \emph{qasam} again
\end{itemize}

\begin{tcolorbox}[breakable,colback=teal!4,colframe=teal!35!black,title={Before you move on}]
What does \hi{क़सम} mean? (\textquotedblleft{}An oath; a vow\textquotedblright{}.) Say it once more.
\end{tcolorbox}
\section[\texorpdfstring{\hi{राज़ी} (rāzī)}{राज़ी (rāzī)}]{\hi{राज़ी} (rāzī) — willing; agreed}

\label{lesson:HI-C543-razi}



\begin{quote}
Before the new one: say the Hindi for a word given, then the Hindi for an oath.
\end{quote}

\subsection*{You'll want to know: \hi{राज़ी}}
\textbf{\hi{राज़ी}} — \emph{rāzī} — \textquotedblleft{}willing; agreed\textquotedblright{}.

\begin{quote}
\hi{पिता} \hi{जी} \hi{मेरी} \hi{बात} \hi{पर} \hi{राज़ी} \hi{हो} \hi{गए।}
\end{quote}

\begin{quote}
\emph{Pitā jī merī bāt par rāzī ho gae.}
\end{quote}

\begin{quote}
Father agreed to what I said.
\end{quote}

\begin{cousinweb}
\textbf{\hi{राज़ी}} (\emph{rāzī}) is usually traced to Persian \emph{rāzī}, \textquotedblleft{}content, consenting\textquotedblright{}.
\end{cousinweb}

\begin{grammarlens}[title={What you've built}]
One more word for saying what happens if.
\end{grammarlens}

\subsection*{Your turn}
\begin{itemize}
\raggedright
  \item \emph{Say it:} \emph{rāzī}
  \item \emph{Say it:} \emph{rāzī}, once more
  \item \emph{Say it:} \emph{Pitā jī merī bāt par rāzī ho gae.} — Father agreed to what I said.
  \item \emph{Recall:} say the Hindi for a word given, then the Hindi for an oath, then say \emph{rāzī} again
\end{itemize}

\begin{tcolorbox}[breakable,colback=teal!4,colframe=teal!35!black,title={Before you move on}]
What does \hi{राज़ी} mean? (\textquotedblleft{}Willing; agreed\textquotedblright{}.) Say it once more.
\end{tcolorbox}
\section[\texorpdfstring{\hi{शपथ} (śapath)}{शपथ (śapath)}]{\hi{शपथ} (śapath) — an oath (formal)}

\label{lesson:HI-C543-sapath}



\begin{quote}
Before the new one: say the Hindi for an oath, then the Hindi for willing.
\end{quote}

\subsection*{You'll want to know: \hi{शपथ}}
\textbf{\hi{शपथ}} — \emph{śapath} — \textquotedblleft{}an oath (formal)\textquotedblright{}.

\begin{quote}
\hi{नए} \hi{मंत्री} \hi{ने} \hi{शपथ} \hi{ली।}
\end{quote}

\begin{quote}
\emph{Nae mantrī ne śapath lī.}
\end{quote}

\begin{quote}
The new minister took the oath.
\end{quote}

Feminine.

A formal word, used for an official oath; everyday speech says \hi{क़सम}.

\begin{cousinweb}
\textbf{\hi{शपथ}} (\emph{śapath}) is a tatsama, taken whole from Sanskrit \emph{śapatha}, \textquotedblleft{}a curse; an oath\textquotedblright{} (Monier-Williams).
\end{cousinweb}

\begin{grammarlens}[title={What you've built}]
One more word for saying what happens if.
\end{grammarlens}

\subsection*{Your turn}
\begin{itemize}
\raggedright
  \item \emph{Say it:} \emph{śapath}
  \item \emph{Say it:} \emph{śapath}, once more
  \item \emph{Say it:} \emph{Nae mantrī ne śapath lī.} — The new minister took the oath.
  \item \emph{Recall:} say the Hindi for an oath, then the Hindi for willing, then say \emph{śapath} again
\end{itemize}

\begin{tcolorbox}[breakable,colback=teal!4,colframe=teal!35!black,title={Before you move on}]
What does \hi{शपथ} mean? (\textquotedblleft{}An oath (formal)\textquotedblright{}.) Say it once more.
\end{tcolorbox}
\section[\texorpdfstring{\hi{पाबंद} (pāband)}{पाबंद (pāband)}]{\hi{पाबंद} (pāband) — bound (by a rule or a time)}

\label{lesson:HI-C543-paband}



\begin{quote}
Before the new one: say the Hindi for willing, then the Hindi for an oath (formal).
\end{quote}

\subsection*{You'll want to know: \hi{पाबंद}}
\textbf{\hi{पाबंद}} — \emph{pāband} — \textquotedblleft{}bound (by a rule or a time)\textquotedblright{}.

\begin{quote}
\hi{वह} \hi{समय} \hi{का} \hi{बहुत} \hi{पाबंद} \hi{है।}
\end{quote}

\begin{quote}
\emph{Vah samay kā bahut pāband hai.}
\end{quote}

\begin{quote}
He always keeps to time.
\end{quote}

\begin{cousinweb}
\textbf{\hi{पाबंद}} (\emph{pāband}) is usually traced to Persian \emph{pāband}, \textquotedblleft{}with feet bound, fettered\textquotedblright{}.
\end{cousinweb}

\begin{grammarlens}[title={What you've built}]
One more word for saying what happens if.
\end{grammarlens}

\subsection*{Your turn}
\begin{itemize}
\raggedright
  \item \emph{Say it:} \emph{pāband}
  \item \emph{Say it:} \emph{pāband}, once more
  \item \emph{Say it:} \emph{Vah samay kā bahut pāband hai.} — He always keeps to time.
  \item \emph{Recall:} say the Hindi for willing, then the Hindi for an oath (formal), then say \emph{pāband} again
\end{itemize}

\begin{tcolorbox}[breakable,colback=teal!4,colframe=teal!35!black,title={Before you move on}]
What does \hi{पाबंद} mean? (\textquotedblleft{}Bound (by a rule or a time)\textquotedblright{}.) Say it once more.
\end{tcolorbox}
